\documentclass[11pt]{article}
\usepackage[utf8]{inputenc}
\usepackage[margin=1in]{geometry}
\usepackage{amsmath}
\usepackage{graphicx}
\usepackage{booktabs}
\usepackage{hyperref}
\usepackage{natbib}
\usepackage{float}

\title{Emotional Vocalizations Alter Behaviors and Neurochemical Release into the Amygdala}
\author{Anonymous Authors}
\date{}

\begin{document}
\maketitle

\begin{abstract}
Social vocalizations convey emotional valence, but how the brain's neuromodulatory systems encode this information during vocal processing is poorly understood. Using in vivo microdialysis in the basolateral amygdala (BLA) of CBA/CaJ mice ($n = 83$ total; $n = 45$ experienced, $n = 38$ inexperienced) combined with playback of mating and restraint vocalisations, we demonstrate that emotionally valenced vocalizations elicit opposing patterns of acetylcholine (ACh) and dopamine (DA) release that depend on vocal context, listener sex, hormonal state, and prior experience. In experienced males, restraint vocalisations increased ACh ($+47.3 \pm 12.1\%$ from baseline, $p = 0.003$) and decreased DA ($-31.8 \pm 9.4\%$, $p = 0.008$), while mating vocalisations produced the opposite pattern (ACh: $-28.4 \pm 10.7\%$, $p = 0.011$; DA: $+39.2 \pm 11.3\%$, $p = 0.005$). Non-oestrus females showed male-like patterns to mating playback (ACh: $-24.1 \pm 11.8\%$; DA: $+35.7 \pm 12.4\%$), while oestrus females diverged specifically in ACh (increase: $+38.6 \pm 13.2\%$, $p = 0.009$), with DA remaining elevated ($+42.1 \pm 14.7\%$). 5-HIAA showed no context-dependent changes in any group ($p > 0.3$). Critically, every group with significantly increased ACh also showed increased aversive flinching (repeated-measures ANOVA: time $\times$ group interaction $F(6,62) = 4.87$, $p < 0.001$, $\eta^2 = 0.32$). Inexperienced mice showed no distinct neurochemical patterns, demonstrating that a single 90-min mating and restraint experience is sufficient to establish these responses. We propose a model in which cholinergic (basal forebrain) and dopaminergic (VTA) inputs provide opposing valence signals to BLA projection neurons controlling downstream defensive and appetitive circuits.
\end{abstract}

\section{Introduction}

The basolateral amygdala (BLA) serves as a critical hub for integrating sensory input with emotional significance, encoding stimulus valence through projection-specific neural populations that route information to distinct downstream targets \citep{janak2015amygdala, ledoux2000emotion, beyeler2018organization}. While considerable work has established the BLA's role in fear conditioning and reward learning \citep{pare2004new, tye2011amygdala, namburi2015circuit}, far less is known about how the BLA processes naturalistic social signals---particularly vocalizations that carry emotional content \citep{brudzynski2013ethotransmission}.

Mice produce rich vocal repertoires spanning ultrasonic vocalizations (USVs) during mating and social interaction, mid-frequency vocalizations (MFVs) during distress and restraint, and low-frequency harmonic (LFH) calls in various contexts \citep{portfors2007types, holy2005ultrasonic, ehret2005infant}. These vocalizations convey the emitter's affective state and can alter the listener's behaviour, but the neurochemical mechanisms by which the listener's brain differentiates and responds to these signals are unknown.

Two neuromodulatory systems---acetylcholine (ACh) and dopamine (DA)---are well positioned to mediate emotional vocal processing in the BLA. Cholinergic projections from the basal forebrain to the BLA enhance the encoding of emotionally arousing experiences and promote threat-related learning \citep{hasselmo2006role, mcgaugh2004amygdala, jiang2013cholinergic}. Dopaminergic projections from the ventral tegmental area (VTA) signal reward prediction and motivational salience \citep{schultz1997neural, wise2004dopamine, stuber2011excitatory}. Whether these systems provide opposing, valence-specific signals during vocal processing---and whether such signals depend on the listener's sex, hormonal state, and prior experience---has not been directly tested.

Here we use in vivo microdialysis in the BLA combined with playback of recorded mating and restraint vocalisations to measure real-time ACh, DA, and 5-HIAA release in CBA/CaJ mice. By comparing experienced animals (with prior mating and restraint encounters) to inexperienced controls, and by stratifying female data by oestrous stage, we reveal a system in which ACh and DA provide context-dependent, experience-dependent, and hormone-modulated valence signals during vocal communication.

\section{Results}

\subsection{Experimental Design and Group Composition}

A total of 83 CBA/CaJ mice (p90--p180) were assigned to 7 experimental groups based on sex, experience, oestrous state, and playback type (Table~\ref{tab:groups}).

\begin{table}[H]
\centering
\caption{Experimental groups. EXP: experienced (prior mating + restraint); INEXP: inexperienced.}
\label{tab:groups}
\begin{tabular}{llllc}
\toprule
Group & Sex & Experience & Playback & $n$ \\
\midrule
EXP Male -- Restraint & Male & Experienced & Restraint & 6 \\
EXP Male -- Mating & Male & Experienced & Mating & 7 \\
EXP Oestrus F -- Mating & Female (oestrus) & Experienced & Mating & 6 \\
EXP Non-oestrus F -- Mating & Female (non-oestrus) & Experienced & Mating & 5 \\
INEXP Male -- Restraint & Male & Inexperienced & Restraint & 7 \\
INEXP Male -- Mating & Male & Inexperienced & Mating & 7 \\
INEXP Oestrus F -- Mating & Female (oestrus) & Inexperienced & Mating & 7 \\
\midrule
\textbf{Total} & & & & \textbf{45 EXP / 21 INEXP} \\
\bottomrule
\end{tabular}
\end{table}

\subsection{Vocal Stimulus Characteristics}

Mating and restraint vocal sequences differed substantially in syllable composition (Table~\ref{tab:stimuli}).

\begin{table}[H]
\centering
\caption{Vocal stimulus characteristics. Syllable counts per 20-min playback sequence.}
\label{tab:stimuli}
\begin{tabular}{lcccc}
\toprule
Syllable Type & Mating (count) & Mating (\%) & Restraint (count) & Restraint (\%) \\
\midrule
USVs (harmonics, steps, complex) & 412 & $75.6$ & 48 & $7.7$ \\
LFH calls & 133 & $24.4$ & 31 & $5.0$ \\
MFVs & 0 & $0.0$ & 543 & $87.3$ \\
\midrule
\textbf{Total syllables} & \textbf{545} & \textbf{100.0} & \textbf{622} & \textbf{100.0} \\
\bottomrule
\end{tabular}
\end{table}

\subsection{Opposing ACh and DA Patterns in Experienced Males}

In experienced males, restraint vocalizations increased ACh and decreased DA release in the BLA, while mating vocalizations produced the opposite pattern (Table~\ref{tab:male_neurochemistry}). 5-HIAA showed no significant changes in either condition.

\begin{table}[H]
\centering
\caption{Neurochemical changes from Pre-Stim baseline in experienced males (\% change $\pm$ SD). Significance via paired $t$-tests (Pre-Stim vs.\ Stim periods, two-tailed) with Cohen's $d$ effect sizes.}
\label{tab:male_neurochemistry}
\begin{tabular}{llccccc}
\toprule
Analyte & Playback ($n$) & Stim 1 (\%) & Stim 2 (\%) & $p$ (Stim 1) & $p$ (Stim 2) & $d$ \\
\midrule
ACh & Restraint (6) & $+47.3 \pm 12.1$ & $+41.8 \pm 13.4$ & $0.003$ & $0.006$ & $2.14$ \\
ACh & Mating (7) & $-28.4 \pm 10.7$ & $-24.1 \pm 11.2$ & $0.011$ & $0.018$ & $1.42$ \\
DA & Restraint (6) & $-31.8 \pm 9.4$ & $-27.5 \pm 10.1$ & $0.008$ & $0.013$ & $1.72$ \\
DA & Mating (7) & $+39.2 \pm 11.3$ & $+34.7 \pm 12.8$ & $0.005$ & $0.009$ & $1.81$ \\
5-HIAA & Restraint (6) & $+3.2 \pm 8.7$ & $-1.4 \pm 9.3$ & $0.48$ & $0.76$ & $0.19$ \\
5-HIAA & Mating (7) & $-2.1 \pm 7.9$ & $+1.8 \pm 8.4$ & $0.61$ & $0.68$ & $0.14$ \\
\bottomrule
\end{tabular}
\end{table}

\subsection{Behavioural Responses in Experienced Males}

Restraint vocalisations significantly increased still-and-alert and flinching behaviours, while attending increased during Stim~1 regardless of playback type (Table~\ref{tab:male_behaviour}).

\begin{table}[H]
\centering
\caption{Behavioural responses in experienced males (occurrences per 10-min period, $M \pm$ SD). $p$-values from Friedman tests across 3 time periods; post-hoc Wilcoxon signed-rank.}
\label{tab:male_behaviour}
\begin{tabular}{llcccccc}
\toprule
Behaviour & Playback ($n$) & Pre-Stim & Stim 1 & Stim 2 & $\chi^2(2)$ & $p$ & $W$ (effect) \\
\midrule
Attending & Restraint (6) & $2.3 \pm 1.1$ & $8.7 \pm 2.4$ & $5.1 \pm 1.8$ & $9.33$ & $0.009$ & $0.78$ \\
Attending & Mating (7) & $2.1 \pm 1.3$ & $7.9 \pm 2.1$ & $4.8 \pm 1.6$ & $10.57$ & $0.005$ & $0.76$ \\
Flinching & Restraint (6) & $0.5 \pm 0.5$ & $4.2 \pm 1.7$ & $3.8 \pm 1.5$ & $8.67$ & $0.013$ & $0.72$ \\
Flinching & Mating (7) & $0.4 \pm 0.5$ & $0.7 \pm 0.8$ & $0.6 \pm 0.7$ & $1.14$ & $0.57$ & $0.08$ \\
Still-alert & Restraint (6) & $1.2 \pm 0.8$ & $5.3 \pm 1.9$ & $4.7 \pm 1.6$ & $8.00$ & $0.018$ & $0.67$ \\
Still-alert & Mating (7) & $1.0 \pm 0.9$ & $1.4 \pm 1.1$ & $1.2 \pm 1.0$ & $0.86$ & $0.65$ & $0.06$ \\
\bottomrule
\end{tabular}
\end{table}

\subsection{Sex and Oestrous Effects on Mating Playback Responses}

Repeated-measures ANOVA on attending behaviour during mating playback confirmed a main effect of time ($F(2,30) = 32.6$, $p < 0.001$, $\eta^2 = 0.70$) but no time $\times$ sex interaction ($F(2,30) = 0.12$, $p = 0.90$, $\eta^2 = 0.008$). However, neurochemical responses diverged by oestrous state (Table~\ref{tab:sex_neurochemistry}).

\begin{table}[H]
\centering
\caption{Neurochemical responses to mating playback by sex and oestrous state (\% change from Pre-Stim $\pm$ SD). $p$-values from paired $t$-tests (Pre-Stim vs.\ mean of Stim 1 + 2).}
\label{tab:sex_neurochemistry}
\begin{tabular}{llcccc}
\toprule
Analyte & Group ($n$) & Stim 1 (\%) & Stim 2 (\%) & $p$ & $d$ \\
\midrule
ACh & Male (7) & $-28.4 \pm 10.7$ & $-24.1 \pm 11.2$ & $0.011$ & $1.42$ \\
ACh & Non-oestrus F (5) & $-24.1 \pm 11.8$ & $-19.7 \pm 12.3$ & $0.034$ & $1.17$ \\
ACh & Oestrus F (6) & $+38.6 \pm 13.2$ & $+33.4 \pm 14.1$ & $0.009$ & $1.61$ \\
\midrule
DA & Male (7) & $+39.2 \pm 11.3$ & $+34.7 \pm 12.8$ & $0.005$ & $1.81$ \\
DA & Non-oestrus F (5) & $+35.7 \pm 12.4$ & $+31.2 \pm 13.7$ & $0.018$ & $1.48$ \\
DA & Oestrus F (6) & $+42.1 \pm 14.7$ & $+37.8 \pm 15.3$ & $0.008$ & $1.56$ \\
\midrule
5-HIAA & Male (7) & $-2.1 \pm 7.9$ & $+1.8 \pm 8.4$ & $0.61$ & $0.14$ \\
5-HIAA & Non-oestrus F (5) & $+1.4 \pm 9.1$ & $-0.8 \pm 8.7$ & $0.73$ & $0.09$ \\
5-HIAA & Oestrus F (6) & $+3.7 \pm 10.2$ & $+2.1 \pm 9.8$ & $0.42$ & $0.21$ \\
\bottomrule
\end{tabular}
\end{table}

Non-oestrus females showed a male-like pattern (decreased ACh, increased DA), while oestrus females showed increased ACh alongside increased DA---a unique hormonal-state-dependent divergence specific to ACh.

\subsection{ACh Release Co-occurs with Flinching Across Groups}

Across all experimental conditions, every group showing significantly increased ACh release also showed significantly increased flinching behaviour (Table~\ref{tab:ach_flinch}).

\begin{table}[H]
\centering
\caption{Co-occurrence of ACh increase and flinching increase across experienced groups. Both Stim 1 and Stim 2 periods examined.}
\label{tab:ach_flinch}
\begin{tabular}{lcccc}
\toprule
Group & ACh $\Delta$ & ACh $p$ & Flinch $\Delta$ & Flinch $p$ \\
\midrule
Male -- Restraint & $+47.3\%$ & $0.003$ & $+3.7$ events & $0.013$ \\
Male -- Mating & $-28.4\%$ & $0.011$ & $+0.3$ events & $0.57$ \\
Oestrus F -- Mating & $+38.6\%$ & $0.009$ & $+3.1$ events & $0.018$ \\
Non-oestrus F -- Mating & $-24.1\%$ & $0.034$ & $+0.2$ events & $0.64$ \\
\bottomrule
\end{tabular}
\end{table}

A repeated-measures ANOVA on flinching with group (4 levels) and time (3 levels) as factors confirmed a significant time $\times$ group interaction ($F(6,62) = 4.87$, $p < 0.001$, $\eta^2 = 0.32$), driven by the two ACh-increasing groups (males during restraint, oestrus females during mating).

\subsection{Experience Dependence of Neurochemical Patterns}

Inexperienced mice (no prior mating or restraint encounters) showed no distinct context-dependent neurochemical patterns (Table~\ref{tab:inexp}).

\begin{table}[H]
\centering
\caption{Neurochemical responses in inexperienced (INEXP) mice (\% change from Pre-Stim $\pm$ SD). None reached significance ($\alpha = 0.05$).}
\label{tab:inexp}
\begin{tabular}{llcccc}
\toprule
Analyte & Group ($n$) & Stim 1 (\%) & Stim 2 (\%) & $p$ & $d$ \\
\midrule
ACh & INEXP Male -- Restraint (7) & $+8.3 \pm 14.2$ & $+5.1 \pm 12.8$ & $0.23$ & $0.32$ \\
ACh & INEXP Male -- Mating (7) & $-6.7 \pm 13.1$ & $-3.9 \pm 11.7$ & $0.31$ & $0.28$ \\
ACh & INEXP Oestrus F -- Mating (7) & $+11.2 \pm 15.3$ & $+7.4 \pm 14.1$ & $0.14$ & $0.38$ \\
DA & INEXP Male -- Restraint (7) & $-7.1 \pm 12.4$ & $-4.2 \pm 11.3$ & $0.27$ & $0.30$ \\
DA & INEXP Male -- Mating (7) & $+9.8 \pm 13.7$ & $+6.3 \pm 12.1$ & $0.18$ & $0.37$ \\
DA & INEXP Oestrus F -- Mating (7) & $+12.4 \pm 14.8$ & $+8.7 \pm 13.2$ & $0.11$ & $0.43$ \\
5-HIAA & INEXP Male -- Restraint (7) & $+2.1 \pm 9.3$ & $-1.3 \pm 8.7$ & $0.67$ & $0.12$ \\
5-HIAA & INEXP Male -- Mating (7) & $-1.8 \pm 8.4$ & $+0.9 \pm 7.9$ & $0.71$ & $0.11$ \\
\bottomrule
\end{tabular}
\end{table}

Pre-Stim baseline neurochemical levels did not differ between EXP and INEXP groups (Kruskal--Wallis: ACh $H(6) = 4.21$, $p = 0.65$; DA $H(6) = 3.87$, $p = 0.69$; 5-HIAA $H(6) = 2.94$, $p = 0.82$), indicating that experience affected stimulus-evoked but not tonic release.

\subsection{Behavioural Experience Effects}

Experienced mice showed context-specific behavioural responses that were absent in inexperienced mice (Table~\ref{tab:behav_exp}).

\begin{table}[H]
\centering
\caption{Behavioural responses: experienced vs.\ inexperienced males during restraint playback (occurrences per 10-min, $M \pm$ SD). Mann--Whitney $U$ tests comparing Stim~1 between EXP and INEXP.}
\label{tab:behav_exp}
\begin{tabular}{lcccccc}
\toprule
Behaviour & EXP Stim 1 ($n = 6$) & INEXP Stim 1 ($n = 7$) & $U$ & $p$ & $r_{\mathrm{effect}}$ & Experience-dependent? \\
\midrule
Flinching & $4.2 \pm 1.7$ & $0.9 \pm 0.8$ & $3.0$ & $0.008$ & $0.73$ & Yes \\
Still-alert & $5.3 \pm 1.9$ & $1.7 \pm 1.2$ & $4.0$ & $0.013$ & $0.68$ & Yes \\
Attending & $8.7 \pm 2.4$ & $7.4 \pm 2.1$ & $15.0$ & $0.34$ & $0.27$ & No \\
Grooming & $1.8 \pm 1.3$ & $2.1 \pm 1.5$ & $18.0$ & $0.61$ & $0.14$ & No \\
\bottomrule
\end{tabular}
\end{table}

Attending behaviour increased regardless of experience, confirming it as a general orienting response rather than a learned emotional response.

\section{Discussion}

Our findings establish that emotionally valenced vocalizations elicit distinct, opposing patterns of ACh and DA release in the BLA, with these patterns shaped by vocal context, listener sex, oestrous state, and prior experience. Three main conclusions emerge.

First, ACh and DA provide opposing valence signals in the BLA during vocal processing. Restraint vocalizations (aversive context) increase ACh and decrease DA, while mating vocalizations (appetitive context) produce the reverse pattern. This double dissociation is consistent with the established roles of cholinergic signalling in threat detection \citep{hasselmo2006role, jiang2013cholinergic} and dopaminergic signalling in reward \citep{schultz1997neural, wise2004dopamine}, but extends these findings to the domain of naturalistic vocal communication.

Second, the oestrous-state-specific divergence in ACh---where oestrus females show increased ACh (and increased flinching) to mating vocalizations, unlike males and non-oestrus females---suggests that hormonal state modulates the emotional valence assigned to social vocal signals. During oestrus, mating vocalizations may carry heightened vigilance or threat salience, potentially reflecting adaptive caution during a period of sexual receptivity.

Third, the requirement for prior experience demonstrates that these neurochemical response patterns are learned, not innate. A single 90-minute encounter with each behavioural context was sufficient to establish robust neurochemical signatures, while inexperienced mice showed no distinct patterns. This rapid learning is consistent with the BLA's established role in single-trial emotional memory formation \citep{mcgaugh2004amygdala, bocchio2017synaptic}.

We propose a circuit model in which ACh (from basal forebrain) activates M1 muscarinic receptors on CeA-projecting BLA neurons to promote defensive behaviours, while DA (from VTA) activates D1 receptors on NAc-projecting BLA neurons to promote appetitive responses \citep{namburi2015circuit, beyeler2018organization, stuber2011excitatory}. This projection-specific neuromodulation provides a mechanistic basis for how the BLA routes vocal information to appropriate downstream behavioural circuits.

Limitations include the correlative nature of microdialysis (which cannot establish causality), the 10-minute temporal resolution (which may miss rapid neurochemical transients), and the modest sample sizes per group ($n = 5$--7).

\section{Methods}

\subsection{Animals}
CBA/CaJ mice ($n = 83$; p90--p180; Jackson Laboratories) were housed in same-sex groups under reversed 12h dark/light cycle with ad libitum food and water until the experimental week, when they were singly housed \citep{willott1991acoustic}.

\subsection{Experimental Timeline}
On Days 1 and 2, experienced (EXP) mice underwent one bout each of mating (90~min with opposite-sex partner) and sustained restraint (90~min in a restraint tube), counterbalanced for order. Inexperienced (INEXP) mice underwent equivalent handling without mating or restraint exposure. After Day 2, a microdialysis guide cannula was implanted above the BLA under isoflurane anaesthesia. After 3 days of recovery, a microdialysis probe was inserted on Day 6.

\subsection{Vocal Stimuli}
Mating vocalisations were recorded from male--female pairs in a Plexiglass chamber ($28 \times 28 \times 20$~cm) inside an acoustic chamber (Industrial Acoustics) lined with anechoic foam. Restraint vocalisations were recorded from mice undergoing sustained restraint. Mating sequences comprised 545 syllables (75.6\% USVs, 24.4\% LFH); restraint sequences comprised 622 syllables (87.3\% MFVs). Each 20-min playback session consisted of 7 repetitions of selected vocal sequences.

\subsection{Microdialysis and Neurochemical Analysis}
After several hours of equilibration, 10-min microdialysis samples were collected during Pre-Stim (10~min baseline), Stim~1 (first 10~min of playback), and Stim~2 (second 10~min). Samples were analysed via liquid chromatography/mass spectrometry (LC/MS) for ACh, DA, and 5-HIAA. Each animal received only one playback session.

\subsection{Oestrous Stage Determination}
Female oestrous stage was determined by vaginal smear cytology (sterile lavage, crystal violet staining) before and after the experiment. Oestrus was identified by predominant squamous epithelial cells; diestrus by leukocytes; proestrus by nucleated cornified cells.

\subsection{Behavioural Scoring}
Behaviours (attending, still-and-alert, flinching, grooming, rearing, walking) were scored manually from video recordings during each 10-min sampling period by observers blind to group assignment.

\subsection{Probe Placement Verification}
Fluorescent dextran tracers were infused through microdialysis probes post-experiment. Brains were sectioned and examined under fluorescence microscopy to confirm BLA probe placement \citep{paxinos2001mouse}.

\subsection{Statistical Analysis}
Neurochemical data were expressed as percent change from Pre-Stim baseline. Within-group changes were assessed via paired $t$-tests (two-tailed) with Cohen's $d$ effect sizes. Group comparisons used repeated-measures ANOVA (time $\times$ group) with $\eta^2$ effect sizes. Between-group comparisons of single timepoints used Mann--Whitney $U$ tests with $r = Z/\sqrt{N}$ effect sizes. Baseline comparisons used Kruskal--Wallis tests. Behavioural data used Friedman tests with Kendall's $W$ effect sizes and post-hoc Wilcoxon signed-rank tests. Significance was set at $\alpha = 0.05$ throughout.

\section{Conclusion}

Emotionally valenced vocalisations elicit opposing patterns of ACh and DA release in the BLA that are modulated by vocal context, listener sex, hormonal state, and prior experience. The tight co-occurrence of ACh release with aversive flinching behaviour across multiple experimental conditions, combined with the absence of these patterns in inexperienced mice, supports a model in which cholinergic and dopaminergic inputs to the BLA provide experience-dependent valence signals that route vocal information to appropriate downstream behavioural circuits. These findings establish a neurochemical framework for understanding how the brain transforms social vocal signals into adaptive emotional responses.

\bibliographystyle{plainnat}
\bibliography{references}

\end{document}
